% Source: https://artofproblemsolving.com/community/c3h1979386_new_problem
% Stable identity: aops-topic-1979386-q01
Let the roots be $r,s,t$. By Vieta's formulas
\[
r+s+t=a,\qquad rs+st+tr=b,\qquad rst=a .
\]
Since $rst=a>0$, no root is $0$, and either all three roots are positive or exactly two are negative.

\emph{Two negative roots.} Write them as $-p,-q$ with $p,q\ge1$ and let $t\ge1$ be the positive root. Then $pqt=a=t-p-q$, so $t(1-pq)=p+q>0$, which forces $pq<1$: impossible.

\emph{All roots positive}, say $1\le r\le s\le t$. Then $rst=r+s+t\le3t$, so $rs\le3$. If $(r,s)=(1,1)$, then $t=t+2$, impossible. If $(r,s)=(1,2)$, then $2t=t+3$, so $t=3$. If $(r,s)=(1,3)$, then $3t=t+4$, so $t=2<s$, contradicting the order.

So the only possible roots are $1,2,3$, which give $a=6$ and $b=1\cdot2+2\cdot3+3\cdot1=11$. For $a=6$ there is exactly one $b$, and for every other $a$ there is none. The answer is $\boxed{6}$.
